% P03 | 1 Introduction | How the work is organised and reproduced
% guide: dev_notes/paragraph_guides/P03_how_the_work_is_organised_and_reproduced.txt
The classes that were written by hand (\texttt{BrownianMotion}, \texttt{symmetric\_random\_walk}, \texttt{CRRBinomialModel}, \texttt{CRRTree} and the two GBM simulators) are used exactly as written.
Every extension is a subclass or a new function placed after them and marked \texttt{[EXT]} in the notebook, so each added result can be traced back to the original code it builds on.
A single notebook produces every figure, table and number in this report, all random numbers are seeded, and the numbers quoted in the text are read from a file generated by the notebook.
A complete run takes about ten minutes on one core.
The scope is deliberately narrow: constant volatility, a constant interest rate and no dividends, except in the one section where a dividend yield is introduced on purpose.
These are the assumptions under which every method in the report can be compared against an exact or near-exact answer.
Full derivations are collected in Appendix~\ref{app:derivations} and the code listing is in Appendix~\ref{app:code}.
